\documentclass[10pt,a4paper]{amsart}
\usepackage{amsmath,amssymb,amsthm,graphicx}
% Retained preamble for the Stacks Project chapter {CH} (shared project preamble).
% The chapter file pinned-original-source.tex only carries \input{{preamble}} on its first line,
% so no preamble text lives inside the chapter itself.  The Stacks Project uses one single shared
% preamble for every chapter: the snapshot stages it three times (sources/stacks-chow/preamble.tex,
% sources/stacks-etale-cohomology/preamble.tex and focused-seeds/AGP/sources/stacks-snapshot-20261011/
% preamble.tex) and those copies are identical in size.  The block below is copied mechanically,
% line for line, from sources/stacks-chow/preamble.tex lines 14-18 and 154-217: the Xy-pic setup and
% the native theorem declarations together with the mathematical macros.  Nothing was added, removed
% or redefined; the wrapper class and the amsmath/amssymb/amsthm/graphicx base packages are supplied
% by the assembler.
% For commutative diagrams we use Xy-pic
% Chapter: stacks-spaces-more-morphisms.tex -- More on Morphisms of Spaces.
\usepackage[all]{xy}

% We use 2cell for 2-commutative diagrams.
\xyoption{2cell}

% Theorem environments.
%
\theoremstyle{plain}
\newtheorem{theorem}[subsection]{Theorem}
\newtheorem{proposition}[subsection]{Proposition}
\newtheorem{lemma}[subsection]{Lemma}

\theoremstyle{definition}
\newtheorem{definition}[subsection]{Definition}
\newtheorem{example}[subsection]{Example}
\newtheorem{exercise}[subsection]{Exercise}
\newtheorem{situation}[subsection]{Situation}

\theoremstyle{remark}
\newtheorem{remark}[subsection]{Remark}
\newtheorem{remarks}[subsection]{Remarks}

\numberwithin{equation}{subsection}

% Macros
%
\def\lim{\mathop{\mathrm{lim}}\nolimits}
\def\colim{\mathop{\mathrm{colim}}\nolimits}
\def\Spec{\mathop{\mathrm{Spec}}}
\def\Hom{\mathop{\mathrm{Hom}}\nolimits}
\def\Ext{\mathop{\mathrm{Ext}}\nolimits}
\def\SheafHom{\mathop{\mathcal{H}\!\mathit{om}}\nolimits}
\def\SheafExt{\mathop{\mathcal{E}\!\mathit{xt}}\nolimits}
\def\Sch{\mathit{Sch}}
\def\Mor{\mathop{\mathrm{Mor}}\nolimits}
\def\Ob{\mathop{\mathrm{Ob}}\nolimits}
\def\Sh{\mathop{\mathit{Sh}}\nolimits}
\def\NL{\mathop{N\!L}\nolimits}
\def\CH{\mathop{\mathrm{CH}}\nolimits}
\def\proetale{{pro\text{-}\acute{e}tale}}
\def\etale{{\acute{e}tale}}
\def\QCoh{\mathit{QCoh}}
\def\Ker{\mathop{\mathrm{Ker}}}
\def\Im{\mathop{\mathrm{Im}}}
\def\Coker{\mathop{\mathrm{Coker}}}
\def\Coim{\mathop{\mathrm{Coim}}}

% Boxtimes
%
\DeclareMathSymbol{\boxtimes}{\mathbin}{AMSa}{"02}

%
% Macros for moduli stacks/spaces
%
\def\QCohstack{\mathcal{QC}\!\mathit{oh}}
\def\Cohstack{\mathcal{C}\!\mathit{oh}}
\def\Spacesstack{\mathcal{S}\!\mathit{paces}}
\def\Quotfunctor{\mathrm{Quot}}
\def\Hilbfunctor{\mathrm{Hilb}}
\def\Curvesstack{\mathcal{C}\!\mathit{urves}}
\def\Polarizedstack{\mathcal{P}\!\mathit{olarized}}
\def\Complexesstack{\mathcal{C}\!\mathit{omplexes}}
% \Pic is the operator that assigns to X its picard group, usage \Pic(X)
% \Picardstack_{X/B} denotes the Picard stack of X over B
% \Picardfunctor_{X/B} denotes the Picard functor of X over B
\def\Pic{\mathop{\mathrm{Pic}}\nolimits}
\def\Picardstack{\mathcal{P}\!\mathit{ic}}
\def\Picardfunctor{\mathrm{Pic}}
\def\Deformationcategory{\mathcal{D}\!\mathit{ef}}
\usepackage{cleveref}
\crefname{originallemma}{original lemma}{original lemmas}
\Crefname{originallemma}{Original lemma}{Original lemmas}
\makeatletter
\newlabel{lemma-resolution-property-in-Noetherian-case}{{0GV0}{source}{}{sourcepointer-1}{}}
\newlabel{lemma-resolution-property-in-Noetherian-case@cref}{{[originallemma][0][]0GV0}{[1][0][]source}{}{}{}}
\makeatother
\crefname{originalreference}{original reference}{original references}
\Crefname{originalreference}{Original reference}{Original references}
\makeatletter
\newlabel{situation-descend-resolution-property}{{0GUY}{source}{}{sourcepointer-2}{}}
\newlabel{situation-descend-resolution-property@cref}{{[originalreference][0][]0GUY}{[1][0][]source}{}{}{}}
\makeatother
\crefname{originalreference}{original reference}{original references}
\Crefname{originalreference}{Original reference}{Original references}
\makeatletter
\newlabel{spaces-cohomology-lemma-finite-presentation-quasi-compact-colimit}{{07U7}{source}{}{sourcepointer-3}{}}
\newlabel{spaces-cohomology-lemma-finite-presentation-quasi-compact-colimit@cref}{{[originalreference][0][]07U7}{[1][0][]source}{}{}{}}
\makeatother
\crefname{originalreference}{original reference}{original references}
\Crefname{originalreference}{Original reference}{Original references}
\makeatletter
\newlabel{spaces-definition-base-change}{{03I5}{source}{}{sourcepointer-4}{}}
\newlabel{spaces-definition-base-change@cref}{{[originalreference][0][]03I5}{[1][0][]source}{}{}{}}
\makeatother
\crefname{originalreference}{original reference}{original references}
\Crefname{originalreference}{Original reference}{Original references}
\makeatletter
\newlabel{spaces-descent-lemma-finite-over-finite-module}{{060Z}{source}{}{sourcepointer-5}{}}
\newlabel{spaces-descent-lemma-finite-over-finite-module@cref}{{[originalreference][0][]060Z}{[1][0][]source}{}{}{}}
\makeatother
\crefname{originalreference}{original reference}{original references}
\Crefname{originalreference}{Original reference}{Original references}
\makeatletter
\newlabel{spaces-limits-lemma-descend-closed-immersion}{{0850}{source}{}{sourcepointer-6}{}}
\newlabel{spaces-limits-lemma-descend-closed-immersion@cref}{{[originalreference][0][]0850}{[1][0][]source}{}{}{}}
\makeatother
\crefname{originalreference}{original reference}{original references}
\Crefname{originalreference}{Original reference}{Original references}
\makeatletter
\newlabel{spaces-limits-lemma-descend-etale}{{07SL}{source}{}{sourcepointer-7}{}}
\newlabel{spaces-limits-lemma-descend-etale@cref}{{[originalreference][0][]07SL}{[1][0][]source}{}{}{}}
\makeatother
\crefname{originalreference}{original reference}{original references}
\Crefname{originalreference}{Original reference}{Original references}
\makeatletter
\newlabel{spaces-limits-lemma-descend-finite}{{084Z}{source}{}{sourcepointer-8}{}}
\newlabel{spaces-limits-lemma-descend-finite@cref}{{[originalreference][0][]084Z}{[1][0][]source}{}{}{}}
\makeatother
\crefname{originalreference}{original reference}{original references}
\Crefname{originalreference}{Original reference}{Original references}
\makeatletter
\newlabel{spaces-limits-lemma-descend-finite-presentation}{{07SK}{source}{}{sourcepointer-9}{}}
\newlabel{spaces-limits-lemma-descend-finite-presentation@cref}{{[originalreference][0][]07SK}{[1][0][]source}{}{}{}}
\makeatother
\crefname{originalreference}{original reference}{original references}
\Crefname{originalreference}{Original reference}{Original references}
\makeatletter
\newlabel{spaces-limits-lemma-descend-invertible-modules}{{0D2X}{source}{}{sourcepointer-10}{}}
\newlabel{spaces-limits-lemma-descend-invertible-modules@cref}{{[originalreference][0][]0D2X}{[1][0][]source}{}{}{}}
\makeatother
\crefname{originalreference}{original reference}{original references}
\Crefname{originalreference}{Original reference}{Original references}
\makeatletter
\newlabel{spaces-limits-lemma-descend-modules-finite-presentation}{{07V7}{source}{}{sourcepointer-11}{}}
\newlabel{spaces-limits-lemma-descend-modules-finite-presentation@cref}{{[originalreference][0][]07V7}{[1][0][]source}{}{}{}}
\makeatother
\crefname{originalreference}{original reference}{original references}
\Crefname{originalreference}{Original reference}{Original references}
\makeatletter
\newlabel{spaces-limits-lemma-finite-in-finite-and-finite-presentation}{{07VR}{source}{}{sourcepointer-12}{}}
\newlabel{spaces-limits-lemma-finite-in-finite-and-finite-presentation@cref}{{[originalreference][0][]07VR}{[1][0][]source}{}{}{}}
\makeatother
\crefname{originalreference}{original reference}{original references}
\Crefname{originalreference}{Original reference}{Original references}
\makeatletter
\newlabel{spaces-limits-lemma-limit-is-affine}{{07SQ}{source}{}{sourcepointer-13}{}}
\newlabel{spaces-limits-lemma-limit-is-affine@cref}{{[originalreference][0][]07SQ}{[1][0][]source}{}{}{}}
\makeatother
\crefname{originalreference}{original reference}{original references}
\Crefname{originalreference}{Original reference}{Original references}
\makeatletter
\newlabel{spaces-limits-proposition-approximate}{{07SU}{source}{}{sourcepointer-14}{}}
\newlabel{spaces-limits-proposition-approximate@cref}{{[originalreference][0][]07SU}{[1][0][]source}{}{}{}}
\makeatother
\crefname{originalreference}{original reference}{original references}
\Crefname{originalreference}{Original reference}{Original references}
\makeatletter
\newlabel{spaces-perfect-definition-resolution-property}{{0GUS}{source}{}{sourcepointer-15}{}}
\newlabel{spaces-perfect-definition-resolution-property@cref}{{[originalreference][0][]0GUS}{[1][0][]source}{}{}{}}
\makeatother
\crefname{originalreference}{original reference}{original references}
\Crefname{originalreference}{Original reference}{Original references}
\makeatletter
\newlabel{spaces-perfect-lemma-resolution-property-goes-up-affine}{{0GUU}{source}{}{sourcepointer-16}{}}
\newlabel{spaces-perfect-lemma-resolution-property-goes-up-affine@cref}{{[originalreference][0][]0GUU}{[1][0][]source}{}{}{}}
\makeatother
\title{The resolution property for one defining sheaf of ideals}
\author{The Stacks Project Authors}
\date{stacks-project:spaces-more-morphisms.tex:local-snapshot-20261011; GNU FDL 1.2 or later}
\begin{document}
\maketitle
\noindent\textit{Editorial note.} Counted span is the author's own proof environment, source lines 13573-13709, immediately following the statement at 13565-13571.\par\medskip
\noindent\textit{Editorial note.} Long proof: the approximation stage cites fourteen cross-chapter pointers (Limits of Spaces, Descent on Spaces, Derived Categories of Spaces); all of them are navigation pointers recorded in assembly.external\_references with their canonical tags.\par\medskip
\noindent\textit{Editorial note.} No step of the author's deduction is outside the span; the Noetherian case is invoked only after the approximation has been constructed inside the span.\par\medskip
\noindent\textit{Allowed standard prerequisites (editorial).}\begin{itemize}
\item The situation of descending the resolution property. chapter situation situation-descend-resolution-property, cited in the statement and used throughout. Fixes the setup: a quasi-compact quasi-separated algebraic space with a finite locally free covering by an affine scheme and the classifying ideal I.
\item Noetherian case of the resolution property. chapter lemma lemma-resolution-property-in-Noetherian-case, cited at proof line 13703. Supplies the conclusion after the descent argument has produced a Noetherian approximation with the same data.
\item Pushforward of a finite module along a finite morphism. Descent on Spaces, lemma spaces-descent-lemma-finite-over-finite-module, cited at proof lines 13577-13578. Keeps finite type quasi-coherent modules finite type under the pushforward along the finite locally free covering.
\item Approximation of an algebraic space by a Noetherian one. Limits of Spaces, proposition spaces-limits-proposition-approximate with the descent lemmas listed in the card, cited from proof line 13690 onward. Realises the space, the covering and the relevant module over a Noetherian base so that the Noetherian case applies.
\item Base change of the resolution property along affine morphisms. Derived Categories of Spaces, lemma spaces-perfect-lemma-resolution-property-goes-up-affine, cited at proof lines 13707-13708. Transfers the property from the Noetherian approximation back to X.
\end{itemize}\medskip
\ifdefined\hypertarget\hypertarget{sourcepointer-1}{}\fi
\noindent\textit{Source reference.} Original lemma 0GV0: sources/stacks-spaces-more-morphisms/pinned-original-source.tex, source lines 13538--13538. Navigation pointer only: the counted proof cites this statement by \textbackslash{}ref as a standing input; the deduction the counted proof has to make is carried out inside the counted span, and no part of that deduction is delegated to this pointer.\par\medskip
\ifdefined\hypertarget\hypertarget{sourcepointer-2}{}\fi
\noindent\textit{Source reference.} Original reference 0GUY: sources/stacks-spaces-more-morphisms/pinned-original-source.tex, source lines 13416--13416. Navigation pointer only: the counted proof cites this statement by \textbackslash{}ref as a standing input; the deduction the counted proof has to make is carried out inside the counted span, and no part of that deduction is delegated to this pointer.\par\medskip
\ifdefined\hypertarget\hypertarget{sourcepointer-3}{}\fi
\noindent\textit{Source reference.} Original reference 07U7: focused-seeds/geometry/S9-spaces-resolve-prop-one-sheaf/ref-tags.txt, source lines 4--4. Navigation pointer only: the counted proof cites this statement by \textbackslash{}ref as a standing input; the deduction the counted proof has to make is carried out inside the counted span, and no part of that deduction is delegated to this pointer.\par\medskip
\ifdefined\hypertarget\hypertarget{sourcepointer-4}{}\fi
\noindent\textit{Source reference.} Original reference 03I5: focused-seeds/geometry/S9-spaces-resolve-prop-one-sheaf/ref-tags.txt, source lines 5--5. Navigation pointer only: the counted proof cites this statement by \textbackslash{}ref as a standing input; the deduction the counted proof has to make is carried out inside the counted span, and no part of that deduction is delegated to this pointer.\par\medskip
\ifdefined\hypertarget\hypertarget{sourcepointer-5}{}\fi
\noindent\textit{Source reference.} Original reference 060Z: focused-seeds/geometry/S9-spaces-resolve-prop-one-sheaf/ref-tags.txt, source lines 6--6. Navigation pointer only: the counted proof cites this statement by \textbackslash{}ref as a standing input; the deduction the counted proof has to make is carried out inside the counted span, and no part of that deduction is delegated to this pointer.\par\medskip
\ifdefined\hypertarget\hypertarget{sourcepointer-6}{}\fi
\noindent\textit{Source reference.} Original reference 0850: focused-seeds/geometry/S9-spaces-resolve-prop-one-sheaf/ref-tags.txt, source lines 7--7. Navigation pointer only: the counted proof cites this statement by \textbackslash{}ref as a standing input; the deduction the counted proof has to make is carried out inside the counted span, and no part of that deduction is delegated to this pointer.\par\medskip
\ifdefined\hypertarget\hypertarget{sourcepointer-7}{}\fi
\noindent\textit{Source reference.} Original reference 07SL: focused-seeds/geometry/S9-spaces-resolve-prop-one-sheaf/ref-tags.txt, source lines 8--8. Navigation pointer only: the counted proof cites this statement by \textbackslash{}ref as a standing input; the deduction the counted proof has to make is carried out inside the counted span, and no part of that deduction is delegated to this pointer.\par\medskip
\ifdefined\hypertarget\hypertarget{sourcepointer-8}{}\fi
\noindent\textit{Source reference.} Original reference 084Z: focused-seeds/geometry/S9-spaces-resolve-prop-one-sheaf/ref-tags.txt, source lines 9--9. Navigation pointer only: the counted proof cites this statement by \textbackslash{}ref as a standing input; the deduction the counted proof has to make is carried out inside the counted span, and no part of that deduction is delegated to this pointer.\par\medskip
\ifdefined\hypertarget\hypertarget{sourcepointer-9}{}\fi
\noindent\textit{Source reference.} Original reference 07SK: focused-seeds/geometry/S9-spaces-resolve-prop-one-sheaf/ref-tags.txt, source lines 10--10. Navigation pointer only: the counted proof cites this statement by \textbackslash{}ref as a standing input; the deduction the counted proof has to make is carried out inside the counted span, and no part of that deduction is delegated to this pointer.\par\medskip
\ifdefined\hypertarget\hypertarget{sourcepointer-10}{}\fi
\noindent\textit{Source reference.} Original reference 0D2X: focused-seeds/geometry/S9-spaces-resolve-prop-one-sheaf/ref-tags.txt, source lines 11--11. Navigation pointer only: the counted proof cites this statement by \textbackslash{}ref as a standing input; the deduction the counted proof has to make is carried out inside the counted span, and no part of that deduction is delegated to this pointer.\par\medskip
\ifdefined\hypertarget\hypertarget{sourcepointer-11}{}\fi
\noindent\textit{Source reference.} Original reference 07V7: focused-seeds/geometry/S9-spaces-resolve-prop-one-sheaf/ref-tags.txt, source lines 12--12. Navigation pointer only: the counted proof cites this statement by \textbackslash{}ref as a standing input; the deduction the counted proof has to make is carried out inside the counted span, and no part of that deduction is delegated to this pointer.\par\medskip
\ifdefined\hypertarget\hypertarget{sourcepointer-12}{}\fi
\noindent\textit{Source reference.} Original reference 07VR: focused-seeds/geometry/S9-spaces-resolve-prop-one-sheaf/ref-tags.txt, source lines 13--13. Navigation pointer only: the counted proof cites this statement by \textbackslash{}ref as a standing input; the deduction the counted proof has to make is carried out inside the counted span, and no part of that deduction is delegated to this pointer.\par\medskip
\ifdefined\hypertarget\hypertarget{sourcepointer-13}{}\fi
\noindent\textit{Source reference.} Original reference 07SQ: focused-seeds/geometry/S9-spaces-resolve-prop-one-sheaf/ref-tags.txt, source lines 14--14. Navigation pointer only: the counted proof cites this statement by \textbackslash{}ref as a standing input; the deduction the counted proof has to make is carried out inside the counted span, and no part of that deduction is delegated to this pointer.\par\medskip
\ifdefined\hypertarget\hypertarget{sourcepointer-14}{}\fi
\noindent\textit{Source reference.} Original reference 07SU: focused-seeds/geometry/S9-spaces-resolve-prop-one-sheaf/ref-tags.txt, source lines 15--15. Navigation pointer only: the counted proof cites this statement by \textbackslash{}ref as a standing input; the deduction the counted proof has to make is carried out inside the counted span, and no part of that deduction is delegated to this pointer.\par\medskip
\ifdefined\hypertarget\hypertarget{sourcepointer-15}{}\fi
\noindent\textit{Source reference.} Original reference 0GUS: focused-seeds/geometry/S9-spaces-resolve-prop-one-sheaf/ref-tags.txt, source lines 16--16. Navigation pointer only: the counted proof cites this statement by \textbackslash{}ref as a standing input; the deduction the counted proof has to make is carried out inside the counted span, and no part of that deduction is delegated to this pointer.\par\medskip
\ifdefined\hypertarget\hypertarget{sourcepointer-16}{}\fi
\noindent\textit{Source reference.} Original reference 0GUU: focused-seeds/geometry/S9-spaces-resolve-prop-one-sheaf/ref-tags.txt, source lines 17--17. Navigation pointer only: the counted proof cites this statement by \textbackslash{}ref as a standing input; the deduction the counted proof has to make is carried out inside the counted span, and no part of that deduction is delegated to this pointer.\par\medskip
\begin{lemma}
\label{lemma-resolution-property-one-sheaf}
In Situation \ref{situation-descend-resolution-property},
the algebraic space $X$ has the resolution property
if and only if $\pi_*\mathcal{I}$ is the quotient of a
finite locally free $\mathcal{O}_X$-module.
\end{lemma}

\begin{proof}
The pushforward $\pi_*\mathcal{G}$ of a finite type
quasi-coherent $\mathcal{O}_Y$-module $\mathcal{G}$ is a finite
type quasi-coherent $\mathcal{O}_X$-module by
Descent on Spaces, Lemma
\ref{spaces-descent-lemma-finite-over-finite-module}.
In particular, if $X$ has the resolution property, then $\pi_*\mathcal{I}$
is the quotient of a finite locally free $\mathcal{O}_X$-module by
Derived Categories of Spaces, Definition
\ref{spaces-perfect-definition-resolution-property}.

\medskip\noindent
Assume that we have a surjection $\mathcal{E} \to \pi_*\mathcal{I}$
for some finite locally free $\mathcal{O}_X$-module $\mathcal{E}$.
In the rest of the proof we show that $X$ has the resolution property
by reducing to the Noetherian case handled in
Lemma \ref{lemma-resolution-property-in-Noetherian-case}.
We suggest the reader skip the rest of the proof.

\medskip\noindent
A first reduction is that we may view $X$ as an algebraic space over
$\Spec(\mathbf{Z})$, see
Spaces, Definition \ref{spaces-definition-base-change}.
(This doesn't affect the conditions nor the conclusion of the lemma.)

\medskip\noindent
By Limits of Spaces, Lemma
\ref{spaces-limits-lemma-finite-in-finite-and-finite-presentation}
we can write $Y = \lim Y_i$ with $Y_i$ finite and of finite
presentation over $X$ and where the transition maps are closed
immersions. Consider the closed subspace $Z = V(\mathcal{I})$ of $Y$.
Since $\mathcal{I}$ is of finite type, the morphism $Z \to Y$ is of
finite presentation. Hence we can find an $i$ and a morphism
$Z_i \to Y_i$ of finite presentation whose base change to $Y$
is $Z \to Y$, see Limits of Spaces, Lemma
\ref{spaces-limits-lemma-descend-finite-presentation}.
For $i' \geq i$ denote $Z_{i'} = Z_i \times_{Y_i} Y_{i'}$.
After increasing $i$ we may assume $Z_i \to Y_i$ is a
closed immersion (of finite presentation), see Limits of Spaces, Lemma
\ref{spaces-limits-lemma-descend-closed-immersion}. Denote
$\mathcal{I}_i \subset \mathcal{O}_{Y_i}$ the ideal
sheaf of $Z_i$ and denote $\pi_i : Y_i \to X$ the structure morphism.
Similarly for $i' \geq i$. Since $Z = \lim_{i' \geq i} Z_{i'}$ we have
$$
\pi_*\mathcal{I} = \colim \pi_{i', *}\mathcal{I}_{i'}
$$
The transition maps in the system are all surjective
as follows from the surjectivity of the maps
$\pi_{i, *}\mathcal{O}_{Y_i} \to \pi_{i', *}\mathcal{O}_{Y_{i'}}$
and the fact that $Z_{i'} = Z_i \times_{Y_i} Y_{i'}$.
By Cohomology of Spaces, Lemma
\ref{spaces-cohomology-lemma-finite-presentation-quasi-compact-colimit}
for some $i' \geq i$ the map $\mathcal{E} \to \pi_*\mathcal{I}$
lifts to a map $\mathcal{E} \to \pi_{i', *}\mathcal{I}_{i'}$.
After increasing $i'$ this map
$\mathcal{E} \to \pi_{i', *}\mathcal{I}_{i'}$
becomes surjective (since if not the colimit of the cokernels,
having surjective transition maps, is nonzero).
This reduces us to the case discussed in the next paragraph.

\medskip\noindent
Assume $X$ is an algebraic space over $\mathbf{Z}$ and that
$Y \to X$ is of finite presentation. By absolute Noetherian
approximation we can write $X = \lim X_i$ as a directed limit,
where each $X_i$ is a quasi-separated algebraic space of finite type
over $\mathbf{Z}$ and the transition morphisms are affine, see
Limits of Spaces, Proposition \ref{spaces-limits-proposition-approximate}.
Since $\pi : Y \to X$ is of finite presentation we can find an $i$
and a morphism $\pi_i : Y_i \to X_i$ of finite presentation
whose base change to $X$ is $\pi$, see
Limits of Spaces, Lemma
\ref{spaces-limits-lemma-descend-finite-presentation}.
After increasing $i$ we may assume $\pi_i$ is finite, see
Limits of Spaces, Lemma
\ref{spaces-limits-lemma-descend-finite}.
Next, we may assume there exists a finite locally
free $\mathcal{O}_{X_i}$-module $\mathcal{E}_i$
whose pullback to $X$ is $\mathcal{E}$, see
Limits of Spaces, Lemma \ref{spaces-limits-lemma-descend-invertible-modules}.
We may also assume there is a map
$\mathcal{E}_i \to \pi_{i, *}\mathcal{O}_{Y_i}$
whose pullback to $X$ is the composition
$\mathcal{E} \to \pi_*\mathcal{I} \to \pi_*\mathcal{O}_Y$, see
Limits of Spaces, Lemma
\ref{spaces-limits-lemma-descend-modules-finite-presentation}.
The cokernel
$$
\mathcal{E}_i \to \pi_{i, *}\mathcal{O}_{Y_i} \to \mathcal{Q}_i \to 0
$$
is a coherent $\mathcal{O}_{Y_i}$-module whose pullback to $X$
is the (finitely presented) cokernel $\mathcal{Q}$ of the map
$\mathcal{E} \to \pi_*\mathcal{O}_Y$. In other words, we have
$\mathcal{Q} = \pi_*(\mathcal{O}_Y/\mathcal{I})$. Consider
the map
$$
\mathcal{E}_i
\otimes_{\mathcal{O}_{X_i}}
\pi_{i, *} \mathcal{O}_{Y_i}
\longrightarrow
\pi_{i, *}\mathcal{O}_{Y_i}
\otimes_{\mathcal{O}_{X_i}}
\pi_{i, *} \mathcal{O}_{Y_i}
\to
\pi_{i, *} \mathcal{O}_{Y_i}
\to
\mathcal{Q}_i
$$
where the second arrow is given by the algebra structure
on $\pi_{i, *}\mathcal{O}_{Y_i}$. The pullback of this map
to $Y$ is zero because the image of $\mathcal{E} \to \pi_*\mathcal{O}_Y$
is the ideal $\pi_*\mathcal{I}$. Hence by
Limits of Spaces, Lemma
\ref{spaces-limits-lemma-descend-modules-finite-presentation}
after increasing $i$ we may assume the displayed composition
is zero. This exactly means that the imag of
$\mathcal{E}_i \to \pi_{i, *}\mathcal{O}_{Y_i}$
is of the form $\pi_{i, *}\mathcal{I}_i$
for some coherent ideal sheaf $\mathcal{I}_i \subset \mathcal{O}_{Y_i}$.
Since $\mathcal{E}_i \to \pi_{i, *}\mathcal{O}_{Y_i}$
pulls back to $\mathcal{E} \to \pi_*\mathcal{O}_Y$ we
see that the pullback of $\mathcal{I}_i$ to $Y$ generates $\mathcal{I}$.
Denote $V_i \subset Y_i$ the open subspace whose complement
is $V(\mathcal{I}_i) \subset Y_i$. Then $V$ is the inverse image
of $V_i$ by the comments above. After increasing $i$ we may assume
that $V_i$ is affine and that $\pi_i|_{V_i} : V_i \to X_i$ is 
\'etale, see
Limits of Spaces, Lemmas
\ref{spaces-limits-lemma-limit-is-affine} and
\ref{spaces-limits-lemma-descend-etale}.
Having said all of this, we may apply
Lemma \ref{lemma-resolution-property-in-Noetherian-case}
to conclude that $X_i$ has the
resolution property. Since $X \to X_i$ is affine
we conclude that $X$ has the resolution property too
by Derived Categories of Spaces, Lemma
\ref{spaces-perfect-lemma-resolution-property-goes-up-affine}.
\end{proof}

\end{document}
